\documentclass[12pt]{article}
\usepackage[utf8]{inputenc}
\usepackage[T1]{fontenc}
\usepackage{geometry}
\geometry{a4paper, margin=1in}
\usepackage{graphicx}
\usepackage{hyperref}
\usepackage{enumitem}
\usepackage{times} % Using a standard serif font
\usepackage{url}

\title{The Hidden Heritage of Haryana: Culture, Manuscripts, and Language}
\author{}
\date{}

\begin{document}

\maketitle

\begin{abstract}
Haryana, often perceived through the lens of its agricultural prosperity and martial traditions, harbours a rich and complex cultural tapestry that remains largely obscured from the broader Indian narrative. This document delves into three foundational pillars of this hidden heritage: its vibrant, ritual-bound folk culture; its precious and neglected manuscript traditions; and its resilient local languages, which form the bedrock of its identity. By exploring these elements, we aim to illuminate the depth and distinctiveness of Haryana's cultural landscape.
\end{abstract}

\section{Vibrant Culture}
The culture of Haryana is a dynamic interplay of ancient pastoral traditions, feudal legacies, and contemporary transformations. Unlike the widely celebrated classical arts of other Indian states, Haryana's cultural expressions are deeply embedded in the rhythms of rural life, agrarian cycles, and oral traditions.

\subsection{Folk Music and Dance}
Central to Haryanvi culture is its folk music, which is predominantly male-dominated and characterized by its powerful, resonant vocals and themes of valor, love, and hardship. The \textbf{Saang} (or \textbf{Swang}) theatre form is a cornerstone, a folk drama that blends music, dialogue, and dance to narrate mythological tales, social issues, and heroic legends like those of \textit{Raja Nal} and the epic \textit{Mahabharata}, which is intrinsically linked to the land of Kurukshetra. The \textbf{Dhol} (a barrel-shaped drum) is the quintessential instrument, providing the rhythmic heartbeat for dances like the \textbf{Khoria} and \textbf{Phag}, performed during festivals like Holi and Teej. The \textbf{Ghoomar} dance, though more famously associated with Rajasthan, has deep roots in the region's communities.

\subsection{Rituals and Fairs}
The culture is also defined by its unique rituals. The \textbf{Gugga Naumi} fair venerates Gugga Pir, a folk deity believed to cure snakebites, showcasing a syncretic blend of Hindu and Islamic folk traditions. The \textbf{Surajkund International Crafts Mela} is a modern manifestation that annually showcases the region's crafts, but the true cultural pulse is felt in village \textbf{jalsas} (gatherings) where wrestlers (\textit{pahalwans}) train in traditional \textit{akharas}, and folk singers perform epic ballads known as \textit{raginis}.

\begin{quote}
\textit{``Haryana's culture is not one of urban sophistication but of rural resilience, where every festival, song, and ritual is tied to the land and the cycles of nature."} (Gupta, 2018)
\end{quote}

\subsection{Cuisine and Attire}
The culinary culture, though simple, is distinct—characterized by dairy products like \textit{bajre ki khichdi}, \textit{lassi}, and \textit{makhan}. The traditional attire, including the \textit{dhoti} and \textit{kesariya pagri} (saffron turban) for men, and the vibrant \textit{ghaghra-choli} for women, remains a potent symbol of identity, especially in rural heartlands.

\section{Manuscripts}
The manuscript wealth of Haryana is a largely uncharted archive, preserved in personal collections, \textit{havelis} (traditional mansions), and under-funded state institutions. These manuscripts are written on palm-leaf, handmade paper, and birch-bark, spanning centuries.

\subsection{Institutional Repositories}
The \textbf{Kurukshetra University} and the \textbf{Haryana Academy of History and Culture} hold significant, yet poorly catalogued, collections. A substantial number of manuscripts pertain to the \textit{Mahabharata}, commentaries on the Vedas, and texts related to the \textit{Panchanga} (Hindu calendar) and astrology, reflecting the region's historical role as a center of Brahminical learning. The \textbf{Chaudhary Charan Singh Haryana Agricultural University} also houses unique \textit{vrikshayurveda} (ancient botany and agriculture) manuscripts, highlighting indigenous agricultural knowledge.

\subsection{Themes and Scripts}
Thematically, the manuscripts can be categorized into:
\begin{itemize}
    \item \textbf{Religious and Philosophical:} Commentaries on the \textit{Gita}, \textit{Puranas}, and works by local saints like \textit{Baba Farid}.
    \item \textbf{Secular and Scientific:} Texts on \textit{jyotisha} (astronomy), \textit{vaidya} (medicine), and \textit{vastu shastra}.
    \item \textbf{Literary:} Poetic works in Braj Bhasha and early Haryanvi, including folk epics and \textit{raginis}.
\end{itemize}

Many of these texts are inscribed in scripts that have become obsolete, such as \textbf{Takri}, \textbf{Landa}, and \textbf{Siddhamatrika}, alongside the more common \textbf{Devanagari} and \textbf{Perso-Arabic} scripts, indicating a confluence of Hindu, Jain, and Sufi literary cultures.

\subsection{Preservation Challenges}
The majority of these manuscripts remain hidden due to neglect, lack of proper conservation infrastructure, and the misconception that Haryana lacks a ``classical'' heritage. Scholars have noted that private family collections are particularly vulnerable to decay and dispersal.

\begin{thebibliography}{9}
\bibitem{sharma1} Sharma, R. K. (2015). \textit{Manuscript Heritage of Haryana: A Survey}. Kurukshetra University Press.
\bibitem{singh2} Singh, J. (2019). ``Unveiling the Lost Archives: Manuscripts in Rural Haryana.'' \textit{Journal of Indian Library and Information Science}, 44(2), 112-125.
\end{thebibliography}

\section{Local Languages of Haryana}
The linguistic landscape of Haryana is far more diverse than the common perception of a single ``Haryanvi'' dialect. It is a continuum of dialects that belong to the Western Hindi and Rajasthani language groups, with significant sociolinguistic dynamics at play.

\subsection{The Dialect Continuum}
The primary language is \textbf{Haryanvi} (also known as \textbf{Bangru} or \textbf{Jatu}), which is the dominant dialect in the central and eastern parts. However, linguistic microcosms exist:
\begin{itemize}
    \item \textbf{Bagri:} Spoken in the western districts like Sirsa and Fatehabad, it shares more features with the Rajasthani language family.
    \item \textbf{Ahirwati:} Spoken in the southern districts (Mahendragarh, Rewari), associated with the Ahir (Yadav) community, it has a distinct phonology and vocabulary.
    \item \textbf{Mewati:} In the far south, the language aligns with the Mewati dialect of Rajasthan.
    \item \textbf{Deswali:} Spoken in the eastern belt bordering Uttar Pradesh.
\end{itemize}

\subsection{Sociolinguistic Status}
Despite being the mother tongue of a majority, Haryanvi and its dialects lack official recognition. Hindi is the official language of the state, leading to a situation of diglossia where the vernacular is stigmatized in formal domains. This has led to a gradual language shift, particularly in urban centers.

\subsection{Linguistic Features}
Haryanvi is characterized by its unique phonetic features: a fast tempo, tonal variations, and the frequent use of the retroflex lateral flap /ɭ/ (often romanized as ``l'' with a diacritic). Its grammar differs from standard Hindi in its postpositions and verb conjugations. For example:
\begin{verbatim}
Standard Hindi: Tum kahan ja rahe ho? (Where are you going?)
Haryanvi: Tu kahan jave se? / Tanne kahan jana hai?
\end{verbatim}

\subsection{Cultural and Literary Erosion}
While Haryanvi has a rich oral literary tradition—including \textit{raginis} (short, metaphorical songs), \textit{phag} (spring songs), and epic ballads like \textit{Bhagat Singh} and \textit{Jaimal Phoola}—its written literary corpus is sparse. The rise of Haryanvi cinema and contemporary folk-pop music has created a recent revival, but scholarly and administrative recognition remains minimal, keeping the full depth of its linguistic heritage hidden from the national mainstream.

\begin{thebibliography}{9}
\bibitem{chaudhary3} Chaudhary, S. (2020). \textit{The Languages of Haryana: A Sociolinguistic Study}. Orient BlackSwan.
\bibitem{kumar4} Kumar, A. (2017). ``Language Endangerment and Vitality: The Case of Haryanvi.'' \textit{Language in India}, 17(5), 23-38.
\bibitem{grierson5} Grierson, G. A. (1904-1928). \textit{Linguistic Survey of India}. Vol. IX, Part I. Office of the Superintendent of Government Printing, India.
\end{thebibliography}

\end{document}